\chapter{哈佛式思考：如何講故事}

本章提煉自 James Wood、Sam Marks、Lauren Groff 與 Nick White 之間的一場座談，由 Samantha Liney-Perfoss 主持，並由 LazyingArt LLC 納入本書。這裡沒有黑板，因此也就沒有可以逐字還原的課堂推導。然而，這場講座確實自有其脊骨。它先為寫作立下一條嚴苛的標準，接著依次談到故事的起源、人物與情節、作家自身的執念與身體、失敗與修改，最後則落在讀者如何檢驗作品是否真正有生命力。我們將緊貼這個次序來梳理。凡是討論中浮現出清楚關係之處，我們都會把它示意地寫出來；但仍須記得，這些公式是對講座的澄清，而不是寫在黑板上的定律。

\section{開場標準：寫作必須付出代價}

這場講座的開頭異常嚴厲。它不是從人物、情節或技藝術語說起，而是先問：究竟什麼樣的寫作，才可能真正影響另一個人？答案是：若文字要在情感上有效，它本身就必須是在情感中寫出來的。這不是一句鼓勵，而是一條標準。

若把它壓縮寫出，開場的主張是
\begin{equation}
\text{對讀者的情感效應} \Rightarrow \text{作家的情感投入}.
\end{equation}
講座立刻又把這一點說得更重。真正重要的，不只是你花了多少時間坐在桌前。有效的寫作，代價不止於時間。它要求的是耗損、暴露與風險。
\begin{equation}
\text{有效的寫作} \Rightarrow C > T,
\end{equation}
其中 \(T\) 是單純花掉的時間，而 \(C\) 則是這份工作更廣義的代價。我們不應把它讀成一條定量不等式。它只是用來說明：真正活著的寫作，要求的遠不只是到場。

同一個觀點，接著又被換成 ``all cylinders firing'' 這個形象來重述。當一件作品真正運轉起來時，作家會感到自己被推到情感與智性能力的邊界上。示意地寫，
\begin{equation}
\text{有效的故事} \Rightarrow E + I,
\end{equation}
其中 \(E\) 代表被充分調動的情感能力，\(I\) 則代表被充分調動的智性能力。同樣地，這只是示意性記號。講座是在說：好的寫作既不只是感受，也不只是思考；它要求兩者同時被激活。

直到這樣強硬的開場之後，主持人才把視野拉寬。說故事是人類生活的一部分。我們熟悉那個作家與頁面搏鬥的形象。當然，技藝是存在的。但講座不希望我們把技藝誤認為故事的全部。因此，真正的問題便在恰到好處的時刻出現了：究竟是什麼，使一個故事令人難忘？

\subsection{Question \& Answer}

\paragraph{Question.} 這場講座對「寫作是否有生命」所提出的第一個檢驗是什麼？

\paragraph{Answer.} 第一個檢驗並不只是潤飾得是否光滑，而是這篇文字是否讓作家付出了超過時間本身的代價，以及情感與智性的努力是否都被充分帶了進來。講座一開始便堅持：頁面上的活性，必須以寫作行動本身的活性為前提。

\section{故事從哪裡開始}

標準一旦立下，主持人便提出第一個實際問題：故事從哪裡開始？幾位作家的回答，都拒絕只承認單一起源。對某位作者來說，種子也許會在心中攜帶多年。對另一位而言，故事在開始寫之前其實根本還不存在。再對另一位來說，一句開場、一個開頭，或一句結尾，都可能先於中段而為整體結構定向。

其中一個特別清楚的說法，是那種長時間醞釀的想法：它原本只是靜止地躺著，直到它與後來某件事件、閱讀或經驗相撞，才忽然獲得足夠的迫切性，變成故事的材料。我們可以把這個關係概括為
\begin{equation}
L + X \Rightarrow U + D + G + W,
\end{equation}
其中 \(L\) 是長久攜帶的想法，\(X\) 是後來發生的碰撞，而 \(U,D,G,W\) 則分別代表迫切性、密度、引力與重量。講座的強調落在這個轉變上：原先只是擁有材料，後來卻變成一種強到足以要求形式的壓力。

第二種說法與其說是矛盾，不如說是互補。故事並不是預先完整地儲存在某處，只等著被抄錄出來。它是在寫作的行動中，才逐漸成為它自己：
\begin{equation}
\text{醞釀} \to \text{第一次寫下} \to \text{展開}.
\end{equation}
這是講座對任何天真靈感觀最重要的修正之一。故事並不總是在寫之前就已完全可知。它往往是在被寫的過程裡，才慢慢變得可讀。

第三種說法引入了「定向」。即使中段還模糊不清，一句開頭、一個起點，或一句結尾，也可能已足夠穩住整個空間：
\begin{equation}
(\text{已知的開頭}) \lor (\text{已知的結尾}) \Rightarrow \text{敘事定向}.
\end{equation}
這一點之所以重要，是因為它解釋了為何幾位截然不同的作家，都可能覺得自己知道 ``該從哪裡開始''，即使這句話對他們各自意味的其實很不一樣。

這一部分最具操作性的引擎，來自人物。有位作者說，寫作往往始於一個特定的人，接著就是兩個立即的問題：這個人想要什麼？又有什麼阻擋了他？
\begin{equation}
\text{故事引擎} = \text{人物} + \text{欲望} + \text{障礙}.
\end{equation}
這確實是一條十分有用的局部公式。它本身就已經包含了壓力，而正是壓力，賦予敘事以運動。

\paragraph{A Worked Example.} 我們可以把座談中幾種關於起源的說法綜合起來，而不假裝它們構成一套普世演算法：
\begin{align}
L &\text{ 可能長時間保持靜止},\\
L + X &\Rightarrow U + D + G + W,\\
U + D + G + W &\Rightarrow \text{第一次寫下},\\
\text{第一次寫下} &\Rightarrow \text{展開}.
\end{align}
這不是一條說故事的定律，而是對講座共通模式的審慎重建：醞釀、激活、首次書寫，以及浮現。

本節最後以一個重要的否定性教訓收束。幾位作家提醒我們，不要幻想可以靠任意的轉折去拯救一份虛弱的草稿。當一個人心裡開始想，``然後這件瘋狂的事就發生了''，結果往往是災難。故事的起點必須讓人感到它是掙得的，而不只是外加的裝飾。

\subsection{Question \& Answer}

\paragraph{Question.} 故事究竟從哪裡開始？

\paragraph{Answer.} 講座給出的答案是複數的。故事可能始於長久的醞釀，始於舊想法與新經驗的碰撞，始於一句開頭，始於一句結尾，也可能始於一個承受壓力的人物。但在幾種說法裡，真正決定性的時刻，都發生在那些材料進入寫作並開始於其中展開之際。

\section{人物、行動與情節光譜}

講座在此把焦點收緊。一旦人物被引入為引擎，下一個問題便是：人物與行動究竟是什麼關係？討論一開始仍然貼近經驗。人物可能會讓作者感到意外，會逐漸顯露自己，也可能在某些場景裡比另一些場景來得更清晰。在戲劇寫作中，人物甚至可能更多是被聽見，而不是被看見。兩個人物一旦被放進同一場景，便像實驗中的元素一樣，你只需觀察會迸出什麼火花。

這種實驗性的語言很重要，因為它直接導向一個理論極點。Aristotle 被召喚進來，作為古典觀點的代表：在那裡，人物從屬於行動：
\begin{equation}
\text{人物} \prec \text{行動}.
\end{equation}
座談並沒有把這一點當作一套完整的現代綱領來接受，而是把它當作場域中的一個極端。當代問題不在於如何絕對服從這條原則，而在於如何去思考那種人物與情節都同時活躍的敘事空間。

這個空間被描述為一條光譜：
\begin{equation}
\text{人物驅動} \leftrightarrow \text{情節驅動}.
\end{equation}
這條光譜的價值，在於診斷。若情節主導得過於凶猛，我們就會得到虛假的推進、任意的扭轉，以及 ``plotting for the sake of plotting''。若人物主導而沒有故事壓力，讀者便可能開始感覺什麼事也沒發生。因此，我們可以把這兩種失敗模式寫成
\begin{equation}
\text{不誠實的硬塞} \;\;\leftrightarrow\;\; \text{沒有故事壓力的人物}.
\end{equation}
講座的核心主張，就是作者必須避開這兩端。

那個滑稽的太空船例子，把這一點說得很尖銳。有時候，一位作者引入某個狂野事件，並不是因為故事真的走到了那裡，而是因為作者想藉此逃離「找出接下來真正會發生什麼」的困難。講座的回應非常細膩：有時候，那艘太空船根本不該出現在第八頁。有時候，它標示的其實是故事真正應該開始的地方。這是一種美麗的修正，因為它把一場拙劣的逃逸，轉化成一個結構問題。

本節同時堅持：人物不是一塊固定不動的內在核心。環境很重要。同樣一對人，在不同地方經歷同樣的分手，不會產生同一場景：
\begin{align}
(\text{同一對情侶},\text{海灘}) &\Rightarrow \text{一種行為場域},\\
(\text{同一對情侶},\text{珠穆朗瑪峰}) &\Rightarrow \text{另一種行為場域}.
\end{align}
因此，人物不只是某個人 ``是'' 什麼；人物也是某個人在情境、場所、壓力與世界之下，會成為什麼。

講座接著提出一句極為醒目的提示：讓你的人物去做那件你自己害怕去做的事。這便是轉軸。到這裡，我們已經從敘事機械學，走向作者與作品之間更深的關係。

\subsection{Question \& Answer}

\paragraph{Question.} 我們應當如何理解人物與情節，而不把其中一方壓扁成另一方？

\paragraph{Answer.} 講座把它們看作彼此都不可或缺的壓力形式。人物帶來欲望、障礙與回應；情節則替這些壓力賦予方向與後果。失敗發生在兩種情況：情節變得任意，或者人物失去了敘事上的張力。好的說故事方式，是讓兩者同時保持活性。

\section{靈魂、執念、地方與身體}

此時，主持人再度把問題放大。如果講座已經談到了人物驅動與情節驅動，那是否也可能存在某種「由作家靈魂所驅動」的寫作？這個問題出現在此處非常恰當，因為前一節其實已經暗示：作者自身的恐懼與勇氣，也許正在人物內部運作。

幾位回答者都把話題引向執念、地方，以及反覆回返的內在提問。有位作者說，自己早期的小說是 author-driven 的，雖然並不是字面上的自傳。另有人說，自己寫出的不同作品，在更深層次上其實是同一個內在問題的不同版本。地方則以異常強烈的方式進入。一個像 Mississippi 這樣的地區，並不只是從外部選來的場景，而是某種嵌進骨頭裡的東西。一個人也許花了很多年想逃離它，然後又花了很多年試圖寫回它之中。

本節最重要的形式性修正是
\begin{equation}
\text{虛構的距離} \not\Rightarrow \text{自我的缺席}.
\end{equation}
一個故事可以與作者的實際生活處境相去甚遠，卻仍然極其私密。這一點之所以重要，是因為講座想保護虛構寫作，不讓它落入兩種相反的誤解：其一，以為個人性必須是字面上的；其二，以為虛構性就必須是非個人的。

通往遙遠材料的橋樑，與其說是傳記性的，不如說是身體性的。感官記憶提供了通道：
\begin{equation}
\text{感官記憶} \Rightarrow \text{進入遙遠虛構世界的通道}.
\end{equation}
講座裡的具體例子極其精準。一個人不必真的住過公社，才能知道手裡捧著一只手工碗、碗裡盛著溫熱燕麥粥時，那種觸感是什麼。這已經是身體知識，而身體知識是可以被轉移的。

在這裡，座談的語言變得更具隱喻性：執念、自我最深處、作家的靈魂、每個人物都是棱鏡式的全像、給讀者多少親近感的滴定。我們應保留這樣的語言，卻不能誤讀它。這個主張並不是一種神祕而空泛的朦朧說法，而是在說：虛構寫作會汲取那些反覆出現的情感與記憶結構，而這些結構即使被翻譯到新的形式、新的時代與新的場景裡，也依然存活。

這裡還有一個重要的實際後果。若一位作者一再回到某些地方、某些恐懼或某些問題，這未必是想像力匱乏的失敗；它反而可能表示，作品已經找到自己真正受壓之處。

\subsection{Question \& Answer}

\paragraph{Question.} 虛構作品如何既能極其個人，又不必是字面上的自傳？

\paragraph{Answer.} 講座的回答是：自我會透過執念、地方、身體記憶、恐懼、欲望與反覆回返的內在問題進入虛構。這些東西不需要靠字面上的自我敘述來完成。作者可以虛構外在情境，卻仍然從自己最深處所能取得的材料出發。

\section{失敗、草稿、修改與讀者}

講座現在來到它對寫作過程最細緻的一節。失敗不再被視作一種寫完之後應該掩飾的羞恥，而是被視作工作中的一種媒介。有位作者公開表示自己熱愛失敗，不是因為失敗本身令人愉快，而是因為它允許遊戲、重新開始，以及不受過早整潔束縛的真正實驗。

因此，講座中的過程關係可以寫成
\begin{equation}
\text{草稿} \to \text{失敗} \to \text{重來} \to \text{故事教會自己}.
\end{equation}
這裡的邏輯十分精確。起草會產生不完美的材料。不完美並不只是被忍受，而是被使用。重來也不只是重複，而是在更清楚的限制下再次嘗試。經過足夠多次之後，故事本身的要求便會開始更清楚地顯露出來。

幾位作家接著比較各自的工作方式。有人從不重讀筆記本，幾乎直到最後都完全停留在紙筆階段。有人在活頁紙、打字稿、列印件與批改之間來回移動。另有人強調，故事會被帶進日常生活之中，因此寫作不只是坐在桌前發生的事，也是在走路、焦慮、試圖解開自己寫出來的結時持續存在的東西。這些方法彼此不同，但講座並不把它們當作互相競爭的教條。它們共同指出的是：過程必須是有意識的，而不同的過程，會讓人與材料形成不同類型的接觸。

其中一個尤其強的洞見，出現在「修改」與「讀者」被連在一起的地方。熱氣騰騰地寫作，會讓糟糕的東西也看起來像天才之作，因為那上面還帶著作者自身的即時性。修改則引進了另一顆心智。因此，一個故事的開頭，不只是作品在時間上的第一部分而已：
\begin{equation}
\text{開頭} \Rightarrow \text{教讀者如何閱讀這個故事}.
\end{equation}
這也就是為什麼開頭總是那麼慢。它不只是入口，更是說明書。

\paragraph{Derivation of the Revision Principle.} 講座的邏輯可以逐步寫成：
\begin{enumerate}
\item 早期起草主要是私人的；暫時而言，作者自己就是唯一的讀者。
\item 修改引入另一位讀者，即使那還只是一位被想像出來的讀者。
\item 開頭的幾段，是那位讀者最先學習作品節奏、尺度與注意方式的地方。
\item 因此，修改會特別重地落在開頭上。
\item 一旦開頭已經正確地教會讀者，後面的運動便能更自由地展開。
\end{enumerate}

講座同時也堅持一種時間上的清醒。在起草的熱度裡，我們可能真誠地相信自己寫出了某種光芒萬丈的東西。隔天一看，當距離恢復，那同一段可能就會呈現完全不同的樣子。這並不是對靈感的背叛，而是判斷終於必要地到來。

\subsection{Question \& Answer}

\paragraph{Question.} 失敗在寫作過程中究竟做了什麼？

\paragraph{Answer.} 失敗會揭露當前嘗試的邊界。一旦它被接受為作品內部的一部分，而不是作品外部的羞辱，它就變得富有生產力。它迫使重來，澄清要求，並逐漸讓故事教會作者，它真正想要的是什麼形狀。

\section{生命力、閱讀、編輯與最後的建議}

最後一個段落開始追問：究竟是什麼使一個故事好？第一個回答是否定式的，而正因如此才富有生產力。故事的好，首先並不取決於它是寫實還是實驗、連續還是碎裂、常規還是怪異。首先決定一切的，是生命力：
\begin{equation}
\text{好的故事} \Rightarrow \text{生命力} = \text{不死} = \text{讓人被吸進去}.
\end{equation}
這就是講座後段的主導標準，而它正好回應開頭那條關於代價的標準。真正曾讓作者付出代價的作品，應當在頁面上產生生命。

講座此時轉向閱讀。閱讀本身被分成兩個階段：
\begin{equation}
\text{蜜月式閱讀} \to \text{七年之癢}.
\end{equation}
第一次閱讀暫停判斷，讓作品直接對我們發生作用。第二次重讀則進入分析，追問某一段為何重要，並開始組裝一套關於作品的論證。這個方法既被當作評論者的方法，也被當作工作坊的方法。

對作者自己而言，生命力也有一個身體性的對應標準。當作品愈來愈接近它真正想成為的樣子時，人是感覺得到的：
\begin{equation}
\text{接近所欲之形} \Rightarrow \text{你能感覺到}.
\end{equation}
講座把這種感受稱為一種震動。這是隱喻性的語言，但並不是空話。它的意思是：某些形式上的成功，首先是靠接觸被知道，然後才靠明確理論被說明。

也就在此刻，講座重新回到開頭的幾句話。我們再次被提醒：在情感上有效的寫作，必須以情感來寫；它付出的代價大於時間本身；它會把作者逼到自己能力所及的邊界。這種回返很重要。它讓圓環閉合。頁面上的生命力，並不與寫出它所付出的代價分離。

下一個實際測試，是把作品朗讀出來。朗讀會把文本從我們身邊拉開，讓我們成為它更冷峻的裁判：
\begin{equation}
\text{朗讀} \Rightarrow \text{疏離} \Rightarrow \text{胡扯偵測器}.
\end{equation}
在默讀時看起來似乎還過得去的句子，一旦在公共聲音中說出來，可能立刻就失敗。講座的意思是：這種失敗不應被辯解掉。它本身就是證據。

這自然帶向那句 ``murder your darlings.'' 但講座再次抗拒把判斷偷換成口號：
\begin{equation}
\text{編輯得好} \neq \text{機械式刪除}.
\end{equation}
有些材料確實是自我縱容，必須被割掉。另一些材料之所以艱難，只是因為它們還沒有被寫得足夠好。因此，講座裡最好的編輯標記，不是 ``把這裡刪掉''，而是 ``這裡寫得更好些''。好的編輯，會辨識生命力薄弱之處，並要求一個更有力的句子，而不是單純要求一份更短的稿子。

最後的建議，把一切重新翻譯回習慣。與作品保持接觸：
\begin{equation}
\text{每天碰觸作品} \Rightarrow \text{與故事維持持續接觸}.
\end{equation}
寫作首先是一種活動，而不只是某種日後才擁有的東西：
\begin{equation}
\text{寫作} = \text{過程} \quad \text{先於} \quad \text{成品}.
\end{equation}
而最後，也不應被外來的規則綁死。講座的收尾幾乎等於是在說：真正重要的規則，最終只有那些由故事自身所生成的規則：
\begin{equation}
\text{規則} = \text{當下這個故事內在的規則}.
\end{equation}
這不是給懶惰開的通行證，而是對注意力提出更高的要求。作者必須把故事聽得足夠仔細，直到它自己的規則變得可聽見為止。

\subsection{Question \& Answer}

\paragraph{Question.} 什麼使一個故事成立？我們又如何知道修改是在幫助它，而不是把它改死？

\paragraph{Answer.} 一個故事之所以成立，是因為它有生命力，因為它活著，並且能把我們吸引進去。當修改提高了這種生命力，不論是藉由去除自我縱容、揭露假音，還是強化那些尚未真正站住的段落，它就是在幫助作品。只有當修改變成機械式地服從規則，而不再重新注意作品最有生命之處時，它才會失敗。

\section{Summary}

這場講座始於代價，終於生命力。介於兩者之間，它展開了一套連貫的說故事觀。故事可能醞釀多年，可能在碰撞中活過來，可能從人物開始，也可能從一句開頭開始，甚至唯有進入寫作行動之中才真正浮現。人物與情節必須取得平衡；自我會透過執念、地方與身體記憶進入虛構；失敗不是作品外部的偶發事故，而是作品媒介的一部分；修改把讀者引進來；而判斷則透過重讀、朗讀與編輯距離抵達。

若把整場講座壓縮成最後一條序列，那便是：
\begin{equation}
\text{壓力} \to \text{寫作} \to \text{失敗} \to \text{修改} \to \text{生命力}.
\end{equation}
這條鏈並不是全部，但它確實是整場討論裡貫穿始終的力線。一個故事起於壓力，在語言裡成形，熬過失敗，從修改中學習，最終則要接受頁面上的生死檢驗。